\documentclass[10pt]{extarticle}
\usepackage[a4paper,margin=1.5cm]{geometry}
\usepackage[T1]{fontenc}
\usepackage{amsmath,amssymb}
\usepackage{graphicx}
\usepackage[export]{adjustbox}
\usepackage{multicol}
\usepackage{enumitem}
\usepackage{booktabs,array,tabularx}
\setlength{\columnsep}{0.7cm}
\setlength{\parindent}{0pt}
\setlength{\parskip}{3pt}
\pagestyle{empty}
\makeatletter
% figure registry (figdefs.tex is generated by the extract step)
\newcommand{\qfigdef}[4]{\@namedef{qf@file@#1}{#2}\@namedef{qf@w@#1}{#3}\@namedef{qf@wide@#1}{#4}}
\newcommand{\qf@check}[1]{\@ifundefined{qf@file@#1}{\PackageError{score2}{Unknown figure id #1}{}}{}}
\newcommand{\qf@img}[4][]{\includegraphics[width=#2\dimexpr\@nameuse{qf@w@#4}pt\relax,max width=#3,max height=0.45\textheight,#1]{\@nameuse{qf@file@#4}}}
% Full-width questions: finalize wraps every question that contains a wide figure in \qfullbegin...\qfullend,
% so the WHOLE question (stem, figures, options) is set across the page in reading order, then the two columns
% resume. The question's head - its number, stem and first figure (\qfullsplit marks the end of the head) -
% is kept on one page, so a number is never left at a page foot with its figure overleaf; the rest may break
% across pages normally, so no big gaps are left. Without a \qfullsplit, the whole question is kept together
% if it is under 45% of a page. Never write these yourself.
\newif\ifqf@full
\newif\ifqf@boxed
\newsavebox{\qf@box}
\newcommand{\qf@open}{\global\qf@fulltrue\par\global\qf@boxedtrue
  \setbox\qf@box\vbox\bgroup\hsize\textwidth\linewidth\textwidth}
\newcommand{\qf@emit}[1]{\par\egroup\global\qf@boxedfalse
  \ifdim\dimexpr\ht\qf@box+\dp\qf@box\relax<#1\textheight\noindent\usebox{\qf@box}\par\else\unvbox\qf@box\fi}
\newcommand{\qfullsplit}{\ifqf@boxed\qf@emit{0.85}\fi}
\newcommand{\qf@close}{\ifqf@boxed\qf@emit{0.45}\fi\global\qf@fullfalse}
\newcommand{\qfullbegin}{\end{multicols}\qf@open}
\newcommand{\qfullend}{\qf@close\begin{multicols}{2}}
\newcommand{\qfullstart}{\qf@open}
\newcommand{\qfulldone}{\qf@close}
% \qfig[scale]{id}: block figure, centred, directly where it is written (size set by the tool).
% Fallback if a wide figure is not inside a full-width question: break out of the columns in place - no floats.
\newcommand{\qfig}[2][1]{\qf@check{#2}%
  \ifqf@full\qf@blk{#1}{#2}\else
  \ifnum\@nameuse{qf@wide@#2}=1 \expandafter\@firstoftwo\else\expandafter\@secondoftwo\fi
  {\end{multicols}\begin{center}\qf@img{#1}{\textwidth}{#2}\end{center}\begin{multicols}{2}}%
  {\qf@blk{#1}{#2}}\fi}
\newcommand{\qf@blk}[2]{\par\smallskip{\centering\qf@img{#1}{\linewidth}{#2}\par}\smallskip}
% \qfigi[scale]{id}: inline figure (inside option lists / tables / minipages); never breaks columns.
\newcommand{\qfigi}[2][1]{\qf@check{#2}\qf@img[valign=c]{#1}{\linewidth}{#2}}
\makeatother
% structure macros - lint reads these to check completeness and figure placement
\newcommand{\papersec}[1]{\par\bigskip{\centering\Large\bfseries PAPER-#1\par}\medskip}
\newcommand{\partsec}[2]{\par\medskip\noindent{\large\bfseries PART-#1: #2}\par\smallskip}
\newcommand{\sect}[1]{\par\smallskip\noindent\textbf{#1}\par}
\newcommand{\qn}[1]{\par\smallskip\noindent\textbf{#1.}~}
\newcommand{\paraq}[2]{\par\smallskip\noindent\textbf{Paragraph for Questions #1 and #2}\par}
\newenvironment{qoptions}{\begin{itemize}[leftmargin=1.8em,itemsep=1pt,topsep=2pt,label={}]}{\end{itemize}}
\newcommand{\opt}[2]{\item[(#1)] #2}
% figure registry (generated by score2_tools.py finalize)
\qfigdef{p1_phy_s0_q01_1}{SRGBT - Test 14 (24 Oct 2021)-img/p1_phy_s0_q01_1.png}{86.3}{0}
\qfigdef{p1_phy_s0_q02_1}{SRGBT - Test 14 (24 Oct 2021)-img/p1_phy_s0_q02_1.png}{137.1}{0}
\qfigdef{p1_phy_s0_q03_1}{SRGBT - Test 14 (24 Oct 2021)-img/p1_phy_s0_q03_1.png}{139.8}{0}
\qfigdef{p1_phy_s0_q04_1}{SRGBT - Test 14 (24 Oct 2021)-img/p1_phy_s0_q04_1.png}{38.5}{0}
\qfigdef{p1_phy_s0_q05_1}{SRGBT - Test 14 (24 Oct 2021)-img/p1_phy_s0_q05_1.png}{121.4}{0}
\qfigdef{p1_phy_s0_q06_1}{SRGBT - Test 14 (24 Oct 2021)-img/p1_phy_s0_q06_1.png}{119.3}{0}
\qfigdef{p1_phy_s0_q07_1}{SRGBT - Test 14 (24 Oct 2021)-img/p1_phy_s0_q07_1.png}{80.0}{0}
\qfigdef{p1_phy_s0_q08_1}{SRGBT - Test 14 (24 Oct 2021)-img/p1_phy_s0_q08_1.png}{78.6}{0}
\qfigdef{p1_phy_s0_q09_1}{SRGBT - Test 14 (24 Oct 2021)-img/p1_phy_s0_q09_1.png}{44.5}{0}
\qfigdef{p1_phy_s0_q10_1}{SRGBT - Test 14 (24 Oct 2021)-img/p1_phy_s0_q10_1.png}{71.2}{0}
\qfigdef{p1_phy_s0_q11_1}{SRGBT - Test 14 (24 Oct 2021)-img/p1_phy_s0_q11_1.png}{41.8}{0}
\qfigdef{p1_phy_s0_q12_1}{SRGBT - Test 14 (24 Oct 2021)-img/p1_phy_s0_q12_1.png}{135.5}{0}
\qfigdef{p1_phy_s0_q12_2}{SRGBT - Test 14 (24 Oct 2021)-img/p1_phy_s0_q12_2.png}{159.6}{0}
\qfigdef{p1_phy_s0_q14_1}{SRGBT - Test 14 (24 Oct 2021)-img/p1_phy_s0_q14_1.png}{114.1}{0}
\qfigdef{p1_phy_s0_q17_1}{SRGBT - Test 14 (24 Oct 2021)-img/p1_phy_s0_q17_1.png}{145.0}{0}
\qfigdef{p1_phy_s0_q19_2}{SRGBT - Test 14 (24 Oct 2021)-img/p1_phy_s0_q19_2.png}{95.3}{0}
\qfigdef{p1_chem_s0_q21_A}{SRGBT - Test 14 (24 Oct 2021)-img/p1_chem_s0_q21_A.png}{101.0}{0}
\qfigdef{p1_chem_s0_q21_B}{SRGBT - Test 14 (24 Oct 2021)-img/p1_chem_s0_q21_B.png}{94.0}{0}
\qfigdef{p1_chem_s0_q21_C}{SRGBT - Test 14 (24 Oct 2021)-img/p1_chem_s0_q21_C.png}{195.0}{0}
\qfigdef{p1_chem_s0_q22_C}{SRGBT - Test 14 (24 Oct 2021)-img/p1_chem_s0_q22_C.png}{50.2}{0}
\qfigdef{p1_chem_s0_q22_D}{SRGBT - Test 14 (24 Oct 2021)-img/p1_chem_s0_q22_D.png}{51.7}{0}
\qfigdef{p1_chem_s0_q23_A}{SRGBT - Test 14 (24 Oct 2021)-img/p1_chem_s0_q23_A.png}{68.6}{0}
\qfigdef{p1_chem_s0_q23_B}{SRGBT - Test 14 (24 Oct 2021)-img/p1_chem_s0_q23_B.png}{181.4}{0}
\qfigdef{p1_chem_s0_q23_C}{SRGBT - Test 14 (24 Oct 2021)-img/p1_chem_s0_q23_C.png}{91.5}{0}
\qfigdef{p1_chem_s0_q23_D}{SRGBT - Test 14 (24 Oct 2021)-img/p1_chem_s0_q23_D.png}{62.3}{0}
\qfigdef{p1_chem_s0_q24_R}{SRGBT - Test 14 (24 Oct 2021)-img/p1_chem_s0_q24_R.png}{63.2}{0}
\qfigdef{p1_chem_s0_q24_A}{SRGBT - Test 14 (24 Oct 2021)-img/p1_chem_s0_q24_A.png}{44.7}{0}
\qfigdef{p1_chem_s0_q24_B}{SRGBT - Test 14 (24 Oct 2021)-img/p1_chem_s0_q24_B.png}{40.4}{0}
\qfigdef{p1_chem_s0_q24_C}{SRGBT - Test 14 (24 Oct 2021)-img/p1_chem_s0_q24_C.png}{40.6}{0}
\qfigdef{p1_chem_s0_q25_A}{SRGBT - Test 14 (24 Oct 2021)-img/p1_chem_s0_q25_A.png}{48.3}{0}
\qfigdef{p1_chem_s0_q25_B}{SRGBT - Test 14 (24 Oct 2021)-img/p1_chem_s0_q25_B.png}{54.0}{0}
\qfigdef{p1_chem_s0_q25_C}{SRGBT - Test 14 (24 Oct 2021)-img/p1_chem_s0_q25_C.png}{55.4}{0}
\qfigdef{p1_chem_s0_q25_D}{SRGBT - Test 14 (24 Oct 2021)-img/p1_chem_s0_q25_D.png}{84.2}{0}
\qfigdef{p1_chem_s0_q26_A}{SRGBT - Test 14 (24 Oct 2021)-img/p1_chem_s0_q26_A.png}{144.3}{0}
\qfigdef{p1_chem_s0_q26_B}{SRGBT - Test 14 (24 Oct 2021)-img/p1_chem_s0_q26_B.png}{206.7}{0}
\qfigdef{p1_chem_s0_q26_C}{SRGBT - Test 14 (24 Oct 2021)-img/p1_chem_s0_q26_C.png}{183.2}{0}
\qfigdef{p1_chem_s0_q26_D}{SRGBT - Test 14 (24 Oct 2021)-img/p1_chem_s0_q26_D.png}{73.7}{0}
\qfigdef{p1_chem_s0_q28_i}{SRGBT - Test 14 (24 Oct 2021)-img/p1_chem_s0_q28_i.png}{35.0}{0}
\qfigdef{p1_chem_s0_q28_ii}{SRGBT - Test 14 (24 Oct 2021)-img/p1_chem_s0_q28_ii.png}{35.3}{0}
\qfigdef{p1_chem_s0_q28_iii}{SRGBT - Test 14 (24 Oct 2021)-img/p1_chem_s0_q28_iii.png}{44.9}{0}
\qfigdef{p1_chem_s0_q28_iv}{SRGBT - Test 14 (24 Oct 2021)-img/p1_chem_s0_q28_iv.png}{47.6}{0}
\qfigdef{p1_chem_s0_q29_A}{SRGBT - Test 14 (24 Oct 2021)-img/p1_chem_s0_q29_A.png}{204.1}{0}
\qfigdef{p1_chem_s0_q29_B}{SRGBT - Test 14 (24 Oct 2021)-img/p1_chem_s0_q29_B.png}{129.0}{0}
\qfigdef{p1_chem_s0_q29_C}{SRGBT - Test 14 (24 Oct 2021)-img/p1_chem_s0_q29_C.png}{136.8}{0}
\qfigdef{p1_chem_s0_q29_D}{SRGBT - Test 14 (24 Oct 2021)-img/p1_chem_s0_q29_D.png}{125.9}{0}
\qfigdef{p1_chem_s0_q30_A}{SRGBT - Test 14 (24 Oct 2021)-img/p1_chem_s0_q30_A.png}{69.1}{0}
\qfigdef{p1_chem_s0_q30_B}{SRGBT - Test 14 (24 Oct 2021)-img/p1_chem_s0_q30_B.png}{115.3}{0}
\qfigdef{p1_chem_s0_q30_C}{SRGBT - Test 14 (24 Oct 2021)-img/p1_chem_s0_q30_C.png}{104.8}{0}
\qfigdef{p1_chem_s0_q30_D}{SRGBT - Test 14 (24 Oct 2021)-img/p1_chem_s0_q30_D.png}{51.8}{0}
\qfigdef{p1_chem_s0_q31_1}{SRGBT - Test 14 (24 Oct 2021)-img/p1_chem_s0_q31_1.png}{83.6}{0}
\qfigdef{p1_chem_s0_q31_A}{SRGBT - Test 14 (24 Oct 2021)-img/p1_chem_s0_q31_A.png}{61.8}{0}
\qfigdef{p1_chem_s0_q31_B}{SRGBT - Test 14 (24 Oct 2021)-img/p1_chem_s0_q31_B.png}{75.0}{0}
\qfigdef{p1_chem_s0_q31_C}{SRGBT - Test 14 (24 Oct 2021)-img/p1_chem_s0_q31_C.png}{74.1}{0}
\qfigdef{p1_chem_s0_q31_D}{SRGBT - Test 14 (24 Oct 2021)-img/p1_chem_s0_q31_D.png}{64.6}{0}
\qfigdef{p1_chem_s0_q32_A}{SRGBT - Test 14 (24 Oct 2021)-img/p1_chem_s0_q32_A.png}{99.4}{0}
\qfigdef{p1_chem_s0_q32_B}{SRGBT - Test 14 (24 Oct 2021)-img/p1_chem_s0_q32_B.png}{119.5}{0}
\qfigdef{p1_chem_s0_q32_C}{SRGBT - Test 14 (24 Oct 2021)-img/p1_chem_s0_q32_C.png}{85.0}{0}
\qfigdef{p1_chem_s0_q32_D}{SRGBT - Test 14 (24 Oct 2021)-img/p1_chem_s0_q32_D.png}{117.4}{0}
\qfigdef{p1_chem_s0_q33_1}{SRGBT - Test 14 (24 Oct 2021)-img/p1_chem_s0_q33_1.png}{302.0}{1}
\qfigdef{p1_chem_s0_q35_A}{SRGBT - Test 14 (24 Oct 2021)-img/p1_chem_s0_q35_A.png}{98.8}{0}
\qfigdef{p1_chem_s0_q37_1}{SRGBT - Test 14 (24 Oct 2021)-img/p1_chem_s0_q37_1.png}{50.4}{0}
\qfigdef{p1_chem_s0_q38_1}{SRGBT - Test 14 (24 Oct 2021)-img/p1_chem_s0_q38_1.png}{50.4}{0}
\qfigdef{p1_chem_s0_q38_2}{SRGBT - Test 14 (24 Oct 2021)-img/p1_chem_s0_q38_2.png}{306.5}{1}
\qfigdef{p1_chem_s0_q40_1}{SRGBT - Test 14 (24 Oct 2021)-img/p1_chem_s0_q40_1.png}{294.7}{1}
\raggedcolumns
\makeatletter\@beginparpenalty=10000\makeatother
\newcolumntype{L}[1]{>{\raggedright\arraybackslash}p{#1}}
\begin{document}
\begin{center}{\Large\bfseries SRGBT --- Test 14 (24 Oct 2021)}\end{center}
\begin{multicols}{2}
\raggedcolumns
\papersec{1}
\partsec{1}{PHYSICS}
\qn{1} A block of mass 10 kg connected to another hollow block of same size and negligible mass, by a spring of spring constant 500 N/m, floats in water as shown in the figure. The compression in the spring is\par
\qfig{p1_phy_s0_q01_1}
($\rho_{\text{water}} = 1\times10^{3}$ kg/m$^3$, g = 10 m/s$^2$)
\begin{qoptions}
\opt{A}{10 cm}
\opt{B}{20 cm}
\opt{C}{50 cm}
\opt{D}{100 cm}
\end{qoptions}
\qn{2} Water is filled in a horizontal cylindrical tank filled with position. The volume of tank is `V' and there is small orifice of area of cross--section `S' at other side of tank at shown in figure. If the `W' is total amount of work done in squeezing all volume of water due to applying a constant force on piston then (Assume time to empty vessel is constant)\par
\qfig{p1_phy_s0_q02_1}
\begin{qoptions}
\opt{A}{W $\propto$ V$^{3/2}$}
\opt{B}{W $\propto$ V$^{3}$}
\opt{C}{W $\propto$ S$^{2}$}
\opt{D}{W $\propto$ V}
\end{qoptions}
\qn{3} During a coating process, a thin, flat tape of width w is pulled through a channel of length L containing a Newtonian fluid of density $\rho$ and dynamic viscosity $\mu$. The fluid is in contact with both sides of the tape. Estimate the force required to pull the tape through the channel if the tape has velocity V and the channel has height H. You may assume that the tape is much thinner than H.\par
\qfig{p1_phy_s0_q03_1}
\begin{qoptions}
\opt{A}{$\dfrac{4\mu \mathrm{VLw}}{\mathrm H}$}
\opt{B}{$\dfrac{2\mu \mathrm{VLw}}{\mathrm H}$}
\opt{C}{$\dfrac{\mu \mathrm{VLw}}{\mathrm H}$}
\opt{D}{$\dfrac{\mu \mathrm{VLw}}{2\mathrm H}$}
\end{qoptions}
\qn{4} A liquid has been poured into a cylindrical vessel of mass M (the mass of the vessel bottom can be ignored) and height H. The linear density of the liquid, that is the ratio of the mass of the liquid column to its height is $\lambda$. The height x of the column of liquid at which the common center of gravity of the liquid plus the vessel is in the lowest position,\par
\qfig{p1_phy_s0_q04_1}
\begin{qoptions}
\opt{A}{H/2}
\opt{B}{M/$\lambda$}
\opt{C}{$\dfrac{\mathrm M}{\lambda}\left[\sqrt{1+\dfrac{\lambda \mathrm H}{\mathrm M}}-1\right]$}
\opt{D}{$\dfrac{\mathrm M}{\lambda\left[\sqrt{1+\dfrac{\lambda \mathrm H}{\mathrm M}}+1\right]}$}
\end{qoptions}
\qn{5} A concrete column is reinforced with an iron rod in the middle. It is subjected to compression by a certain load. Assuming that the modulus of elasticity of concrete is one-tenth that of iron and the area of cross section of the iron is one-twentieth that of concrete, find the portion of the load acting on the concrete.\par
\qfig{p1_phy_s0_q05_1}
\begin{qoptions}
\opt{A}{$\left(\dfrac23\right)^{\text{rd}}$ of load}
\opt{B}{$\left(\dfrac13\right)^{\text{rd}}$ of load}
\opt{C}{$\left(\dfrac14\right)^{\text{th}}$ of load}
\opt{D}{$\left(\dfrac25\right)^{\text{th}}$ of load}
\end{qoptions}
\qn{6} A solid cone of density d has $\dfrac13$ rd of its length in a fluid of density d$_1$, other $\dfrac13$ rd in air and the remaining $\dfrac13$ rd in a fluid of density d$_2$, as shown in figure. Given that d$_1 = \dfrac35$ d: the cone is in equilibrium in position shown with the contact force from the upper container's bottom being negligibly small, atmospheric pressure due to air is p$_0$ and force by atmospheric pressure on the middle part in air is F. Pick correct option :\par
\qfig{p1_phy_s0_q06_1}
\begin{qoptions}
\opt{A}{d$_2$ = 10 d}
\opt{B}{F = P$_0\times 2\pi$h$^2$}
\opt{C}{F = P$_0\times 3\pi$h$^2$}
\opt{D}{d$_2$ = 20d}
\end{qoptions}
\qn{7} A wedge having length, breadth and height each equal to `a' moves in medium-1 having pressure p$_0$. It then enters to medium-2 having pressure 2p$_0$. Assuming all the surfaces to be smooth, the change in kinetic energy of the wedge when it completely enters medium-2 is\par
\qfig{p1_phy_s0_q07_1}
\begin{qoptions}
\opt{A}{$-$p$_0$a$^3$}
\opt{B}{p$_0$a$^3$}
\opt{C}{$-$p$_0$(a$^3$/2)}
\opt{D}{$-$p$_0$(a$^3$/3)}
\end{qoptions}
\qn{8} A cone of height h and base of radius r is submerged in a fluid as shown in figure such that half of its height is outside the fluid. The angle of contact between the fluid and cone is $\theta$ and surface tension for the fluid is T. The net force experienced by cone due to surface tension is\par
\qfig{p1_phy_s0_q08_1}
\begin{qoptions}
\opt{A}{$\pi r\mathrm T\cos\theta$}
\opt{B}{$\pi r\mathrm T\sin\theta$}
\opt{C}{$\pi r\mathrm T\cos\left[\theta-\tan^{-1}\left(\dfrac rh\right)\right]$}
\opt{D}{$\pi r\mathrm T\cos\left[\theta-\tan^{-1}\left(\dfrac hr\right)\right]$}
\end{qoptions}
\qn{9} A soap film (surface tension = T) is formed on a rectangular wire frame as shown in figure. A small, light, inextensible loop of thread of length ``$\ell$'' is gently placed over the soap film. The loop stays on the film in the irregular fashion in which it was placed. Now, a hole is pricked in the film inside the loop with a needle and the thread finally comes to equilibrium position. Choose the correct option(s)\par
\qfig{p1_phy_s0_q09_1}
\begin{qoptions}
\opt{A}{The thread will take a square shape}
\opt{B}{The thread will take a circular shape}
\opt{C}{Tension in the thread will be 2T$\ell$}
\opt{D}{Tension in the thread will be $\dfrac{\mathrm T\ell}{\pi}$}
\end{qoptions}
\qn{10} If the velocity profile is given by $v=\dfrac23 y-y^2$ (m/sec.) y is in meter above the bed ($\eta$ = 0.863 Ns/m$^2$) (Total depth of the liquid is 1/3 m.)\par
\qfig{p1_phy_s0_q10_1}
\begin{qoptions}
\opt{A}{Shear stress at $y$ = 0.15 m is 0.63 N/m$^2$}
\opt{B}{Shear stress at $y$ = 0.15 m is 0.316 N/m$^2$}
\opt{C}{Maximum shear stress is 0.575 N/m$^2$}
\opt{D}{Maximum shear stress is 1.15 N/m$^2$}
\end{qoptions}
\qn{11} Linear mass density of the given vertical rod varies as $\lambda=\lambda_0 x$ where $\lambda_0$ is a constant and x is distance from upper end. Youngs modulus and cross sectional area of the rod is Y and A respectively.\par
\qfig{p1_phy_s0_q11_1}
\begin{qoptions}
\opt{A}{elongation of rod due to its own weight is $\dfrac{\lambda_0 gL^3}{3YA}$}
\opt{B}{elongation of rod due to its own weight is $\dfrac{2\lambda_0 gL^3}{3YA}$}
\opt{C}{Tension as a function of x is T = $\lambda_0 g\dfrac{(L^2-x^2)}{3}$}
\opt{D}{Tension as a function of x is T = $\lambda_0 g\dfrac{(L^2-x^2)}{2}$}
\end{qoptions}
\qn{12} Which of the following is/are \textbf{CORRECT} about the following system (take $\pi=\sqrt{10}$ )\par
\qfig{p1_phy_s0_q12_1}
\begin{qoptions}
\opt{A}{h$_\mathrm A$ should be 16 m for water column to be remains steady in container B}
\opt{B}{h$_\mathrm A$ should be 32 m for water column to be remains steady in container B}
\opt{C}{Q should be nearly 4.5 litre/sec for water column to be remains steady in both the containers}
\opt{D}{Q should be nearly 18 litre/sec for water column to be remains steady in both the containers}
\end{qoptions}
\paraq{13}{14} Consider a situation in which there is a Newtonian fluid between two long coaxial uniformly moving cylinders of radius R and 2R. (Use Poisuille's assumptions)\par
\qfig{p1_phy_s0_q12_2}
\qn{13} Force per unit length req to keep inner cylinder moving. Take $\eta$ as coefficient of viscosity
\begin{qoptions}
\opt{A}{$4\pi\eta \mathrm v_0$}
\opt{B}{$2\pi\eta \mathrm v_0$}
\opt{C}{$\pi\eta \mathrm v_0$}
\opt{D}{None of these}
\end{qoptions}
\qn{14} At what radius speed of fluid is zero.
\begin{qoptions}
\opt{A}{$\sqrt2$ R}
\opt{B}{$\dfrac32$R}
\opt{C}{$\dfrac{5\mathrm R}{4}$}
\opt{D}{None of these}
\end{qoptions}
\paraq{15}{16} A container of large uniform cross-section area A resting on a horizontal surface, holds two immiscible, non-viscous and incompressible liquids of densities d and 2d each of height H/2 as shown in figure. The lower density liquid is open to the atmosphere having pressure P$_0$.\par
\qfig{p1_phy_s0_q14_1}
\qn{15} A homogeneous solid cylinder of length $l$($l$ $<$ H/2) cross sectional area A/5 is immersed such that it floats with its axis vertical at the interface of two liquids with length $\dfrac{\ell}{4}$ in the denser liquid. Find the density of solid cylinder.
\begin{qoptions}
\opt{A}{$\dfrac{3d}{4}$}
\opt{B}{$\dfrac32 d$}
\opt{C}{$\dfrac{5d}{4}$}
\opt{D}{$\dfrac{7d}{4}$}
\end{qoptions}
\qn{16} The cylinder is removed and original arrangement is restored. A tiny hole of area (S$<<$A) is punched on the vertical side of the container at a height h $\left(\mathrm h<\dfrac{\mathrm H}{2}\right)$. from the bottom of the container, find the initial speed of efflux of the liquid at the hole.
\begin{qoptions}
\opt{A}{$\sqrt{\dfrac g2}(3\mathrm H-4\mathrm h)$}
\opt{B}{$\sqrt{g(3\mathrm H-4\mathrm h)}$}
\opt{C}{$\sqrt{g(4\mathrm H-3\mathrm h)}$}
\opt{D}{$\sqrt{\dfrac g2}(\mathrm H-\mathrm h)$}
\end{qoptions}
\qn{17} A soap bubble is blown at the end of a capillary tube of radius a and length L. When the other end is left open, the bubble begins to deflate. The radius of the bubble as a function of time is given by $\mathrm R_0\left[1-\dfrac{\mathrm a^4\mathrm Tt}{\mathrm m\eta \mathrm L\mathrm R_0^4}\right]^{\frac1{\mathrm n}}$ (where initial radius of the bubble was R$_0$. Surface tension of soap solution is T. It is known that volume flow rate through a tube of radius a and length L is given by Poiseuille's equation - Q $=\dfrac{\pi \mathrm a^4\Delta \mathrm P}{8\eta \mathrm L}$ ; $\Delta$P is pressure difference at the two ends of the tube and $\eta$ is coefficient of viscosity. Assume that the bubble remains spherical) Find the value of m $\times$ n.\par
\qfig{p1_phy_s0_q17_1}
\qn{18} An ice cube of size $\ell$ = 20 cm is floating in a tank (base area A = 50 cm $\times$ 50 cm) partially filled with water. Density of water is $\rho_{\text{w}}$ = 1000 kg/m$^3$ and that of ice is $\rho_{\text{i}}$ = 900 kg/m$^3$. The decrease in gravitational potential energy when ice melts completely is found to be K times 0.08 J. Find K. (g = 10 m/s$^2$)
\qn{19} The bar in figure below is made from a single piece of material. It consists of two segments of equal length $\dfrac{\mathrm L_0}{2}$ and cross-sectional area A and 2A. The elastic potential energy stored under the action of an axial force F is given in multiples of $\dfrac{\mathrm F^2\mathrm L_0}{\mathrm{AY}}$ as $\dfrac{\mathrm n}{24}$. What is the value of n ?\par
\qfig{p1_phy_s0_q19_2}
\qn{20} A rod of length L is suspended from the ceiling by one of its ends. The density of the material of the rod is $\rho$ and its Young's modulus is Y. The relative increment of its volume is given by $\dfrac{\rho \mathrm{gL}}{\mathrm{nY}}(1-\mathrm n\sigma)$. What is the value of n? Given $\sigma$ is the Poisson's ratio.
\raggedcolumns
\papersec{1}
\partsec{2}{CHEMISTRY}
\qn{21} Select incorrect option.
\begin{qoptions}
\opt{A}{\qfigi{p1_chem_s0_q21_A} is correct order of stability.}
\opt{B}{\qfigi{p1_chem_s0_q21_B} is correct order of dipole moment.}
\opt{C}{\qfigi{p1_chem_s0_q21_C} is correct order of heat of combustion.}
\opt{D}{Bond order of C--C bond in conjugate base of cyclo penta--1,3--diene is 7/5.}
\end{qoptions}
\qn{22} Which of the following is less acidic than benzoic acid.
\begin{qoptions}
\opt{A}{HCOOH}
\opt{B}{PhSO$_3$H}
\opt{C}{\qfigi{p1_chem_s0_q22_C}}
\opt{D}{\qfigi{p1_chem_s0_q22_D}}
\end{qoptions}
\qn{23} Select correct statement.
\begin{qoptions}
\opt{A}{\qfigi{p1_chem_s0_q23_A} is correct order of reactivity for EAS reaction.}
\opt{B}{\qfigi{p1_chem_s0_q23_B} is correct order of pK$_\mathrm b$.}
\opt{C}{\qfigi{p1_chem_s0_q23_C} is correct order of resonance energy.}
\opt{D}{\qfigi{p1_chem_s0_q23_D} is correct order of dipole moment.}
\end{qoptions}
\qn{24} Which of the following product is possible in given reaction.\par
\qfigi{p1_chem_s0_q24_R}
\begin{qoptions}
\opt{A}{\qfigi{p1_chem_s0_q24_A}}
\opt{B}{\qfigi{p1_chem_s0_q24_B}}
\opt{C}{\qfigi{p1_chem_s0_q24_C}}
\opt{D}{All are possible}
\end{qoptions}
\qn{25} Select correct IUPAC name.
\begin{qoptions}
\opt{A}{\qfigi{p1_chem_s0_q25_A} $\Rightarrow$ Diethanoyl chloride}
\opt{B}{\qfigi{p1_chem_s0_q25_B} $\Rightarrow$ Ethoxy carbonyl ethanoic acid.}
\opt{C}{\qfigi{p1_chem_s0_q25_C} $\Rightarrow$ Methyl--2--methyl butanoate}
\opt{D}{\qfigi{p1_chem_s0_q25_D}\par $\Rightarrow$ 2,6--dichloro--4,5--diethenylocta--2,7--diene}
\end{qoptions}
\qn{26} Which of the following reactions is incorrect ?
\begin{qoptions}
\opt{A}{\qfigi{p1_chem_s0_q26_A}}
\opt{B}{\qfigi{p1_chem_s0_q26_B}}
\opt{C}{\qfigi{p1_chem_s0_q26_C}}
\opt{D}{Styrene \qfigi{p1_chem_s0_q26_D}}
\end{qoptions}
\qn{27} Which of the following can't liberate CO$_2$ gas on reaction with NaHCO$_3$.
\begin{qoptions}
\opt{A}{Squaric acid}
\opt{B}{Barbituric acid}
\opt{C}{Carbolic acid}
\opt{D}{Benzene sulphonic acid}
\end{qoptions}
\qn{28} (i) \qfigi{p1_chem_s0_q28_i} (ii) \qfigi{p1_chem_s0_q28_ii} (iii) \qfigi{p1_chem_s0_q28_iii} (iv) \qfigi{p1_chem_s0_q28_iv}\par
Correct order of acidic nature is:
\begin{qoptions}
\opt{A}{iii $>$ i $>$ ii $>$ iv}
\opt{B}{i $>$ iii $>$ ii $>$ iv}
\opt{C}{ii $>$ iii $>$ i $>$ iv}
\opt{D}{ii $>$ iii $>$ iv $>$ i}
\end{qoptions}
\qn{29} Which of the following reactions is/are correct ?
\begin{qoptions}
\opt{A}{\qfigi{p1_chem_s0_q29_A}}
\opt{B}{\qfigi{p1_chem_s0_q29_B}}
\opt{C}{\qfigi{p1_chem_s0_q29_C}}
\opt{D}{\qfigi{p1_chem_s0_q29_D}}
\end{qoptions}
\qn{30} Which of the following is/are correct for enol content ?
\begin{qoptions}
\opt{A}{\qfigi{p1_chem_s0_q30_A}}
\opt{B}{\qfigi{p1_chem_s0_q30_B}}
\opt{C}{\qfigi{p1_chem_s0_q30_C}}
\opt{D}{\qfigi{p1_chem_s0_q30_D}}
\end{qoptions}
\qn{31} Which of the following is/are conformational of (p) ?\par
\qfigi{p1_chem_s0_q31_1}
\begin{qoptions}
\opt{A}{\qfigi{p1_chem_s0_q31_A}}
\opt{B}{\qfigi{p1_chem_s0_q31_B}}
\opt{C}{\qfigi{p1_chem_s0_q31_C}}
\opt{D}{\qfigi{p1_chem_s0_q31_D}}
\end{qoptions}
\qn{32} Select incorrect relations.
\begin{qoptions}
\opt{A}{\qfigi{p1_chem_s0_q32_A} $\Rightarrow$ Positional isomer}
\opt{B}{\qfigi{p1_chem_s0_q32_B} $\Rightarrow$ Identical}
\opt{C}{\qfigi{p1_chem_s0_q32_C} $\Rightarrow$ Conformational isomers}
\opt{D}{\qfigi{p1_chem_s0_q32_D} $\Rightarrow$ Chain isomer}
\end{qoptions}
\qfullbegin
\paraq{33}{34} Consider the following molecules having same molecular formula.\par
\qfig{p1_chem_s0_q33_1}
\qfullend
\qn{33} Which isomerism is not observed in between above compounds.
\begin{qoptions}
\opt{A}{Chain Isomer}
\opt{B}{Position Isomer}
\opt{C}{Metamer Isomer}
\opt{D}{Functional Isomer}
\end{qoptions}
\qn{34} Most acidic in the above molecules.
\begin{qoptions}
\opt{A}{d}
\opt{B}{c}
\opt{C}{f}
\opt{D}{e}
\end{qoptions}
\paraq{35}{36} In organic compounds electrons of $\sigma$ \& $\pi$--bonds are show displacement from their normal position due to various reason. It is known as electronic displacement effect. For electronic displacement effect give answer of following questions.
\qn{35} For inductive effect, select incorrect statement.
\begin{qoptions}
\opt{A}{\qfigi{p1_chem_s0_q35_A}, inductive effect of $-$Me on $-$COOH is same in both compounds.}
\opt{B}{Inductive effect is distance dependent effect.}
\opt{C}{Magnitude of +I effect does not depend upon geometry of compound.}
\opt{D}{Maximum cationic group can show $-$I effect but it is not necessary for all.}
\end{qoptions}
\qn{36} Select correct statement for mesomeric effect.
\begin{qoptions}
\opt{A}{All catonic groups can show $-$M effect.}
\opt{B}{All catonic groups can show $+$M effect.}
\opt{C}{$\dot{\mathrm C}\mathrm H_2-\mathrm{CN}$ is less stable than $\dot{\mathrm C}\mathrm H_3$ due to $-$M effect of $-$CN group.}
\opt{D}{Some groups can show both $+$M and $-$M depending upon condition.}
\end{qoptions}
\qn{37} How many resonating structures are possible for \qfigi{p1_chem_s0_q37_1} including this structure ?
\qfullbegin
\qn{38} Give sum of serial number of following compound which are more reactive than \qfigi{p1_chem_s0_q38_1} towards EAS.\par
\qfig{p1_chem_s0_q38_2}
\qfullend
\qn{39} How many conformers are possible for butane. (Consider rotation about C$_2$--C$_3$ bond)
\qfullbegin
\qn{40} Consider the following compound :\par
\qfig{p1_chem_s0_q40_1}
\qfullsplit
Serial number of most basic compound is (P)\par
Serial number of most acidic compound is (Q)\par
Total number of aromatic compound = (R)\par
Number of compound is represented in bracket.\par
Find the value of P + Q + R.
\raggedcolumns
\qfullend
\papersec{1}
\partsec{3}{MATHEMATICS}
\qn{41} Consider an ellipse E : $\dfrac{x^2}{4}+\dfrac{y^2}{3}=1$ , S is its focus and line L is its corresponding directrix. Curve `C' is the locus of a variable point which moves on xy plane such that its distance from S and line L are equal. Number of common tangent(s) of curve C and E is-
\begin{qoptions}
\opt{A}{0}
\opt{B}{1}
\opt{C}{2}
\opt{D}{3}
\end{qoptions}
\qn{42} Let S$_1$ is a hyperbola and S$_2$ be an ellipse whose centres are at origin. The transverse axis of S$_1$ and major axis of S$_2$ lie along the x-axis and y-axis respectively. Let H be the hyperbola $x^2-y^2=7$. The straight line 4x $-$ 3y $=$ 7. touches the curves H, S$_1$ and S$_2$ at P, Q and R respectively, Suppose. PQ$=\dfrac52=$QR If e$_1$ and e$_2$ are eccentricities of S$_1$ and S$_2$ respectively. The which of the following is true :
\begin{qoptions}
\opt{A}{e$_1=\dfrac{23}{15}$}
\opt{B}{e$_2=\sqrt{\dfrac12}$}
\opt{C}{15e$_1{}^2$+4e$_2{}^2$ = 24}
\opt{D}{e$_1$ is any integer}
\end{qoptions}
\qn{43} A tangent to the hyperbola $\dfrac{x^2}{4}-\dfrac{y^2}{1}=1$ meets ellipse $x^2+4y^2=4$ at two distinct points. The locus of mid points of this chord is :
\begin{qoptions}
\opt{A}{$(x^2+4y^2)^2=4(x^2-4y^2)$}
\opt{B}{$(x^2-4y^2)=4(x^2+4y^2)$}
\opt{C}{$(x^2-4y^2)^2=4(4x^2+y^2)$}
\opt{D}{$(x^2+4y^2)^2=4(4x^2-y^2)$}
\end{qoptions}
\qn{44} Four common tangents of ellipse $\dfrac{x^2}{4}+\dfrac{y^2}{1}=1$ and hyperbola $\dfrac{x^2}{5}-\dfrac{y^2}{3}=1$ make a parallelogram whose area is
\begin{qoptions}
\opt{A}{$4\sqrt{17}$}
\opt{B}{$\dfrac{17}{2}$}
\opt{C}{17}
\opt{D}{34}
\end{qoptions}
\qn{45} If the length of the major axis of an ellipse is 100 and eccentricity $=\dfrac12$, then the area of the rectangle formed by joining the vertices of latus recta of the ellipse is equal to :
\begin{qoptions}
\opt{A}{2500}
\opt{B}{3750}
\opt{C}{5000}
\opt{D}{6250}
\end{qoptions}
\qn{46} A tangent of the ellipse $\dfrac{x^2}{a^2}+\dfrac{y^2}{b^2}=1$ is normal to the hyperbola $\dfrac{x^2}{4}-\dfrac{y^2}{1}=1$ and it has equal intercepts with positive x and y axes, then the value of $a^2+b^2$ is
\begin{qoptions}
\opt{A}{5}
\opt{B}{$\dfrac{25}{3}$}
\opt{C}{16}
\opt{D}{$\dfrac{25}{9}$}
\end{qoptions}
\qn{47} If A(1, 2) and B(9, $-$6) are two points on the parabola $y^2=4x$. The ellipse $\dfrac{(x-1)^2}{\lambda}+\dfrac{y^2}{1}=1$; $\lambda>0$ passing through A and B intersect the parabola again at C and D. The point of intersection of normals at C and D lies at
\begin{qoptions}
\opt{A}{2x $-$ y $-$ 12 = 0}
\opt{B}{x $-$ 2y $-$ 10 = 0}
\opt{C}{x $+$ 2y $-$ 12 = 0}
\opt{D}{x $-$ 2y $-$ 12 = 0}
\end{qoptions}
\qn{48} The point on the hyperbola $\dfrac{x^2}{24}-\dfrac{y^2}{18}=1$ which is nearest to the line 3x + 2y+1=0 is
\begin{qoptions}
\opt{A}{(6, 3)}
\opt{B}{($-$6, 3)}
\opt{C}{(6, $-$3)}
\opt{D}{($-$6, $-$3)}
\end{qoptions}
\qn{49} If a tangent of slope $\dfrac13$ of the ellipse $\dfrac{x^2}{a^2}+\dfrac{y^2}{b^2}=1$ (a $>$ b) is normal to the circle $x^2+y^2+2x+2y+1=0$ then:
\begin{qoptions}
\opt{A}{Maximum value of ab is $\dfrac23$}
\opt{B}{$a\in\left(\sqrt{\dfrac25},2\right)$}
\opt{C}{$a\in\left(\dfrac23,2\right)$}
\opt{D}{Maximum value of ab is 1}
\end{qoptions}
\qn{50} If $a^2(x^2+y^2-2x-4y+5)=(5x-12y+3)^2$ is an ellipse, then the value of `a' can be -
\begin{qoptions}
\opt{A}{$\sqrt{168}$}
\opt{B}{17}
\opt{C}{21}
\opt{D}{$\sqrt{170}$}
\end{qoptions}
\qn{51} If a variable chord of the hyperbola $x^2-y^2=9$ touches the parabola $y^2=12x$, then locus of middle points of these chords is expressed as $x^3+\lambda_1xy^2+\lambda_2y^2=0$ ($\lambda_1,\lambda_2\in$ I), then-
\begin{qoptions}
\opt{A}{$\lambda_1+\lambda_2=2$}
\opt{B}{$\lambda_1^2+\lambda_2^2=10$}
\opt{C}{Number of divisors of $\left(\lambda_1^2+\lambda_2^2\right)$ is 4}
\opt{D}{Number of divisors of $\left(\lambda_1^2+\lambda_2^2\right)$ is 6}
\end{qoptions}
\qn{52} If (a sec$\theta$, b tan$\theta$) and (a sec$\phi$, b tan$\phi$) are the ends of a focal chord of $\dfrac{x^2}{a^2}-\dfrac{y^2}{b^2}=1$, then $\tan\dfrac\theta2\tan\dfrac\phi2$ equals to
\begin{qoptions}
\opt{A}{$\dfrac{e-1}{e+1}$}
\opt{B}{$\dfrac{1-e}{1+e}$}
\opt{C}{$\dfrac{1+e}{1-e}$}
\opt{D}{$\dfrac{e+1}{e-1}$}
\end{qoptions}
\paraq{53}{54} Consider the curves C$_1$ : $y=\dfrac{8}{27}x^3$ and C$_2$ : $y=(x+a)^2$.
\qn{53} The range of `a' for which there exists two common tangents to the curves C$_1$ and C$_2$ other than x-axis is :
\begin{qoptions}
\opt{A}{$(0,\infty)$}
\opt{B}{$\left(-\dfrac12,\infty\right)-\{0\}$}
\opt{C}{$\left(-\dfrac32,\infty\right)-\{0\}$}
\opt{D}{$[1,\infty)$}
\end{qoptions}
\qn{54} If a = 4 ,then the area bounded by the two tangent lines mentioned in the previous question and the given parabola C$_2$ is :
\begin{qoptions}
\opt{A}{100}
\opt{B}{144}
\opt{C}{180}
\opt{D}{200}
\end{qoptions}
\paraq{55}{56} If a circle touches the parabola $y^2=4x$ at A(9, 6) and is tangents to x-axis then
\qn{55} The maximum possible radius of circle is $p+q\sqrt r$ , where p, q, r $\in$ even integers then p+q+r is :
\begin{qoptions}
\opt{A}{88}
\opt{B}{78}
\opt{C}{68}
\opt{D}{58}
\end{qoptions}
\qn{56} The $y$-intercept of common normal through A is $\alpha$ times x-intercept of common normal through A and the area of triangle formed by tangent, normal at A to parabola and x-axis is $\beta$ square units, then $\alpha+\beta$ is :
\begin{qoptions}
\opt{A}{68}
\opt{B}{65}
\opt{C}{63}
\opt{D}{61}
\end{qoptions}
\qn{57} Let A$_1$ and A$_2$ are the vertices of the conic C$_1$ : $4(x-3)^2+9(y-2)^2-36=0$ and a point P is moving in the plane such that $|\mathrm{PA}_1-\mathrm{PA}_2|=3\sqrt2$ , then locus of P is another conic C$_2$. If D$_1$ denotes distance between foci of conic C$_2$. D$_2$ denotes product of the perpendiculars from the points A$_1$,A$_2$ upon any tangent drawn to conic C$_2$ and D$_3$ denotes length of the tangent drawn from any point on auxiliary circle of conic C$_1$ to the auxiliary circle of the conic C$_2$, then $\left(\dfrac{\mathrm D_1\mathrm D_2}{\mathrm D_3{}^2}\right)$ is equal to
\qn{58} If a normal chord of the parabola $2y^2=x$ has slope $(2-\sqrt3)$ , then the area of triangle formed by normal chord of the parabola and the tangents drawn at the extremities of the normal chord is :
\qn{59} If the locus of mid points of portions of tangents intercepted between axes of hyperbola $\dfrac{x^2}{16}-\dfrac{y^2}{9}=1$ is $\dfrac{\mathrm A}{x^2}-\dfrac9{y^2}=\mathrm B$, then $\left(\dfrac{\mathrm A+\mathrm B}{4}\right)$ is equal to :
\qn{60} The tangent at point P on the hyperbola $\dfrac{x^2}{a^2}-\dfrac{y^2}{b^2}=1$ passes through the point (0, $-$b) and the normal at point P passes through the point $\left(2\sqrt2\,a,0\right)$. If e denote the eccentricity of hyperbola then find the value of e$^2$.
\end{multicols}
\end{document}
